\chapter{Observability and Diagnostics}

Operational visibility is essential for production systems. This chapter specifies the observability model for gert: structured logging, metrics, distributed tracing, health probes, and diagnostics tooling. It is distinct from \S~12 (Evidence, Tracing, and Resumption): \textbf{§12 specifies immutable audit records for compliance and resumption}; this chapter specifies \textbf{operational signals for operators} to answer questions like \textit{is the system healthy?}, \textit{what is slow?}, \textit{where did this run fail?}, and \textit{what is the resource utilization?}

The design follows the three-pillar observability model~\cite{google-sre-book} adapted for workflow execution: metrics for aggregated performance, distributed tracing for execution flow, and structured logging for detailed context. All three are instrumented using industry standards (OpenTelemetry~\cite{otel-spec}, W3C Trace Context~\cite{w3c-trace-context}, RFC 3339 timestamps~\cite{rfc3339}) to ensure interoperability with observability platforms.

% ================================================================
\section{Observability Model}
% ================================================================

Gert applies the three-pillar observability model to workflow execution. Each pillar serves a distinct purpose:

\subsection{Metrics: Aggregated Performance Indicators}

\textbf{Purpose:} Track aggregate performance, throughput, error rates, and resource utilization over time.

\textbf{Use Cases:}
\begin{itemize}
  \item \textit{Alerting:} Trigger alerts when error rates exceed thresholds or latency degrades
  \item \textit{Capacity Planning:} Measure run throughput and active concurrency to inform scaling decisions
  \item \textit{SLO Monitoring:} Track percentile latencies (p50, p95, p99) for step execution times
\end{itemize}

\textbf{Metric Types:}
\begin{itemize}
  \item \textbf{Counter:} Monotonically increasing (total runs, total errors)
  \item \textbf{Gauge:} Point-in-time measurement (active runs, loaded extensions)
  \item \textbf{Histogram:} Distribution of values (run duration, tool invocation latency)
\end{itemize}

\textbf{Exposition Format:} Prometheus OpenMetrics format exposed via HTTP endpoint (\texttt{GET /metrics}) when running in server mode (\texttt{gert serve --http}).

\subsection{Distributed Tracing: Execution Flow Visualization}

\textbf{Purpose:} Capture the hierarchical execution tree with timing and outcome for every run, step, and tool invocation.

\textbf{Use Cases:}
\begin{itemize}
  \item \textit{Root-Cause Analysis:} Identify which step in a failed run caused the failure
  \item \textit{Latency Diagnosis:} Find which tool invocations are slow
  \item \textit{Cross-System Correlation:} Link gert spans to HTTP API call traces via W3C Trace Context propagation
\end{itemize}

\textbf{Implementation:} OpenTelemetry SDK with OTLP exporter~\cite{otel-spec}. Spans emitted for every run (root span), step (child spans), and tool invocation (grandchild spans).

\textbf{Span Hierarchy:}
\begin{verbatim}
gert.run (run_id=...)
├─ gert.step (step_id=backup-db, step_type=cli)
│  ├─ gert.tool.invoke (tool=ssh, exit_code=0)
│  └─ gert.governance.check (policy=allow-ssh)
├─ gert.step (step_id=deploy, step_type=tool)
│  └─ gert.tool.invoke (tool=kubectl, exit_code=0)
└─ gert.step (step_id=verify, step_type=manual)
   └─ gert.input.resolve (provider=approval-gate)
\end{verbatim}

\textbf{Trace Backends:} Jaeger~\cite{jaeger-docs}, Zipkin~\cite{zipkin-docs}, Datadog, New Relic, AWS X-Ray (all OTLP-compatible).

\subsection{Structured Logging: Detailed Contextual Events}

\textbf{Purpose:} Emit detailed, machine-parseable log lines with run and step context for debugging and incident investigation.

\textbf{Use Cases:}
\begin{itemize}
  \item \textit{Debugging:} Search logs by \texttt{run\_id} to see all activity for a specific run
  \item \textit{Error Investigation:} View full error messages and stack traces in context
  \item \textit{Compliance Audit:} Query logs for all runs by a specific actor or runbook
\end{itemize}

\textbf{Log Format:} JSON, one object per line, emitted to \texttt{stderr} by default.

\textbf{Required Fields:}
\begin{itemize}
  \item \texttt{timestamp} (RFC 3339~\cite{rfc3339})
  \item \texttt{level} (debug, info, warn, error)
  \item \texttt{msg} (human-readable message)
  \item \texttt{run\_id} (when in run context)
  \item \texttt{step\_id} (when in step context)
\end{itemize}

\textbf{Aggregation:} Loki, Elasticsearch, Splunk, CloudWatch Logs (any system that can parse JSON logs).

% ================================================================
\section{OpenTelemetry Integration}
% ================================================================

Gert integrates OpenTelemetry as the standard tracing layer. All spans, metrics, and logs follow the OpenTelemetry specification~\cite{otel-spec} to ensure portability across observability backends.

\subsection{Span Hierarchy and Attributes}

OpenTelemetry spans form a hierarchical tree representing the execution structure. Each span level has required attributes for filtering and analysis.

\subsubsection{Root Span: \texttt{gert.run}}

The root span represents the entire run from start to completion.

\textbf{Span Name:} \texttt{gert.run}

\textbf{Required Attributes:}
\begin{minted}{text}
runbook.id       = "deploy-app-v3"      // runbook identifier
runbook.version  = "1.2.3"              // version (if specified)
actor.id         = "alice@example.com"  // human or service principal
run.id           = "20260418T143022-a3f9c1"  // unique run ID
run.mode         = "real"               // "real" | "replay" | "dry-run"
client.type      = "cli"                // "cli" | "tui" | "web" | "mcp"
\end{minted}

\textbf{Span Status:} 
\begin{itemize}
  \item \texttt{OK} if run outcome is \texttt{success}
  \item \texttt{ERROR} if run outcome is \texttt{failure}, \texttt{cancelled}, or \texttt{governance-blocked}
\end{itemize}

\subsubsection{Child Span: \texttt{gert.step}}

Each step execution produces a child span under \texttt{gert.run}.

\textbf{Span Name:} \texttt{gert.step}

\textbf{Required Attributes:}
\begin{minted}{text}
step.id          = "backup-db"          // step identifier from YAML
step.type        = "cli"                // "cli" | "manual" | "tool" | "invoke"
step.index       = 2                    // zero-based position in sequence
run.id           = "20260418T143022-a3f9c1"  // run correlation
step.name        = "Backup production database"  // human-readable name
\end{minted}

\textbf{Optional Attributes:}
\begin{minted}{text}
step.outcome     = "success"            // "success" | "failure" | "skipped"
step.retries     = 2                    // number of retry attempts
\end{minted}

\textbf{Span Events:} Emit span events for key transitions:
\begin{itemize}
  \item \texttt{governance.check.started} — governance evaluation started
  \item \texttt{governance.check.passed} — governance check passed
  \item \texttt{approval.requested} — human approval requested
  \item \texttt{approval.granted} — approval decision logged
\end{itemize}

\subsubsection{Grandchild Span: \texttt{gert.tool.invoke}}

Tool invocations create grandchild spans under \texttt{gert.step}.

\textbf{Span Name:} \texttt{gert.tool.invoke}

\textbf{Required Attributes:}
\begin{minted}{text}
tool.name        = "kubectl"            // tool identifier
tool.transport   = "exec"               // "exec" | "ssh" | "http"
tool.exit_code   = 0                    // process exit code (exec only)
run.id           = "20260418T143022-a3f9c1"  // run correlation
step.id          = "deploy"             // parent step
\end{minted}

\textbf{Optional Attributes:}
\begin{minted}{text}
tool.args        = ["apply", "-f", "deployment.yaml"]  // redacted args
tool.duration_ms = 1234                 // invocation latency
\end{minted}

\subsubsection{Grandchild Span: \texttt{gert.input.resolve}}

Input provider resolutions create grandchild spans.

\textbf{Span Name:} \texttt{gert.input.resolve}

\textbf{Required Attributes:}
\begin{minted}{text}
input.name       = "db_password"        // input key from runbook
provider.type    = "vault"              // provider type
provider.address = "https://vault:8200" // endpoint (if applicable)
run.id           = "20260418T143022-a3f9c1"
\end{minted}

\subsubsection{Grandchild Span: \texttt{gert.governance.check}}

Governance evaluations create grandchild spans.

\textbf{Span Name:} \texttt{gert.governance.check}

\textbf{Required Attributes:}
\begin{minted}{text}
policy.name      = "allow-ssh"          // policy identifier
policy.outcome   = "allow"              // "allow" | "deny"
run.id           = "20260418T143022-a3f9c1"
step.id          = "backup-db"
\end{minted}

\subsection{Span Status and Error Handling}

OpenTelemetry spans have three status codes:
\begin{itemize}
  \item \texttt{UNSET} — default, no explicit status set
  \item \texttt{OK} — operation completed successfully
  \item \texttt{ERROR} — operation failed
\end{itemize}

\textbf{Status Mapping:}
\begin{itemize}
  \item \texttt{gert.run} span:
    \begin{itemize}
      \item \texttt{OK} if run outcome is \texttt{success}
      \item \texttt{ERROR} if run outcome is \texttt{failure}, \texttt{cancelled}, \texttt{governance-blocked}
    \end{itemize}
  \item \texttt{gert.step} span:
    \begin{itemize}
      \item \texttt{OK} if step outcome is \texttt{success} or \texttt{skipped}
      \item \texttt{ERROR} if step outcome is \texttt{failure}
    \end{itemize}
  \item \texttt{gert.tool.invoke} span:
    \begin{itemize}
      \item \texttt{OK} if \texttt{exit\_code == 0}
      \item \texttt{ERROR} if \texttt{exit\_code != 0} or invocation crashes
    \end{itemize}
\end{itemize}

When a span is marked \texttt{ERROR}, the runtime MUST record the error message as a span event or exception:

\begin{minted}{go}
span.RecordError(err, trace.WithAttributes(
  attribute.String("error.type", "tool_invocation_failed"),
  attribute.String("error.message", err.Error()),
))
\end{minted}

\subsection{Context Propagation}

Gert propagates trace context across system boundaries using W3C Trace Context~\cite{w3c-trace-context}.

\subsubsection{W3C Trace Context Headers}

When a tool invokes an HTTP API, the runtime MUST inject W3C headers:

\begin{verbatim}
traceparent: 00-4bf92f3577b34da6a3ce929d0e0e4736-00f067aa0ba902b7-01
tracestate: gert=run_id:20260418T143022-a3f9c1
\end{verbatim}

The \texttt{traceparent} header contains:
\begin{itemize}
  \item \texttt{version} (00)
  \item \texttt{trace-id} (128-bit, globally unique)
  \item \texttt{parent-id} (64-bit span ID)
  \item \texttt{trace-flags} (sampling decision)
\end{itemize}

The \texttt{tracestate} header carries vendor-specific metadata. Gert includes \texttt{run\_id} as baggage for cross-system correlation.

\subsubsection{Baggage}

The \texttt{run\_id} is propagated as OpenTelemetry baggage to all child spans and external systems:

\begin{minted}{go}
baggage.Set("run_id", runID)
\end{minted}

This allows downstream systems to tag their own spans with the originating gert run ID, enabling end-to-end tracing across multi-system workflows.

\subsection{OTLP Exporter Configuration}

The runtime supports two trace export modes:

\subsubsection{OTLP Export (Production)}

Configure via environment variable or CLI flag:

\begin{minted}{bash}
export OTEL_EXPORTER_OTLP_ENDPOINT=http://jaeger:4317
gert run deploy.yaml
\end{minted}

Or:

\begin{minted}{bash}
gert run deploy.yaml --otel-endpoint=http://jaeger:4317
\end{minted}

The runtime uses the OTLP/gRPC exporter by default. For HTTP, use:

\begin{minted}{bash}
export OTEL_EXPORTER_OTLP_ENDPOINT=http://jaeger:4318
export OTEL_EXPORTER_OTLP_PROTOCOL=http/protobuf
\end{minted}

\subsubsection{Console Export (Development)}

For local development without a trace backend:

\begin{minted}{bash}
gert run deploy.yaml --otel-console
\end{minted}

This prints spans to \texttt{stderr} in a human-readable format.

\subsection{Sampling Strategy}

Gert uses \textbf{parent-based sampling} with a configurable default rate:

\begin{itemize}
  \item If an incoming request already has a trace context (e.g., gert invoked via HTTP from another traced system), \textbf{inherit the parent's sampling decision}.
  \item If no parent context exists, sample at the configured rate (default: 100\% = always sample).
\end{itemize}

Configuration:

\begin{minted}{bash}
export OTEL_TRACES_SAMPLER=parentbased_traceidratio
export OTEL_TRACES_SAMPLER_ARG=0.1  # sample 10% of traces
\end{minted}

For production systems with high throughput, consider \texttt{0.01} (1\%) or \texttt{0.001} (0.1\%) to reduce backend load.

% ================================================================
\section{Metrics Catalog}
% ================================================================

Gert exposes the following metrics in Prometheus OpenMetrics format. All metrics are prefixed with \texttt{gert\_}.

\subsection{Run Metrics}

\begin{table}[h]
\centering
\begin{tabular}{|p{4.5cm}|p{1.5cm}|p{3.5cm}|p{5.5cm}|}
\hline
\textbf{Metric} & \textbf{Type} & \textbf{Labels} & \textbf{Description} \\
\hline
\texttt{gert\_runs\_total} & Counter & \texttt{outcome}, \texttt{runbook\_id} & Total run completions by outcome (\texttt{success}, \texttt{failure}, \texttt{cancelled}) \\
\hline
\texttt{gert\_run\_duration\_seconds} & Histogram & \texttt{outcome}, \texttt{runbook\_id} & End-to-end run duration (start to completion) \\
\hline
\texttt{gert\_active\_runs} & Gauge & — & Currently executing runs \\
\hline
\end{tabular}
\caption{Run-level metrics}
\end{table}

\textbf{Histogram Buckets:} \texttt{gert\_run\_duration\_seconds} uses exponential buckets: \texttt{0.5, 1, 2, 5, 10, 30, 60, 120, 300, 600, 1800, 3600} (seconds).

\subsection{Step Metrics}

\begin{table}[h]
\centering
\begin{tabular}{|p{4.5cm}|p{1.5cm}|p{3.5cm}|p{5.5cm}|}
\hline
\textbf{Metric} & \textbf{Type} & \textbf{Labels} & \textbf{Description} \\
\hline
\texttt{gert\_steps\_total} & Counter & \texttt{outcome}, \texttt{step\_type} & Total step completions by outcome and type \\
\hline
\texttt{gert\_step\_duration\_seconds} & Histogram & \texttt{step\_type} & Per-step duration by type (\texttt{cli}, \texttt{manual}, \texttt{tool}, \texttt{invoke}) \\
\hline
\end{tabular}
\caption{Step-level metrics}
\end{table}

\textbf{Histogram Buckets:} \texttt{gert\_step\_duration\_seconds} uses buckets: \texttt{0.1, 0.5, 1, 2, 5, 10, 30, 60, 120, 300} (seconds).

\subsection{Tool Invocation Metrics}

\begin{table}[h]
\centering
\begin{tabular}{|p{4.5cm}|p{1.5cm}|p{3.5cm}|p{5.5cm}|}
\hline
\textbf{Metric} & \textbf{Type} & \textbf{Labels} & \textbf{Description} \\
\hline
\texttt{gert\_tool\_invocations\_total} & Counter & \texttt{tool\_name}, \texttt{outcome} & Tool invocation count by tool and outcome (\texttt{success}, \texttt{failure}) \\
\hline
\texttt{gert\_tool\_duration\_seconds} & Histogram & \texttt{tool\_name} & Tool invocation latency \\
\hline
\end{tabular}
\caption{Tool invocation metrics}
\end{table}

\textbf{Histogram Buckets:} \texttt{gert\_tool\_duration\_seconds} uses buckets: \texttt{0.01, 0.05, 0.1, 0.5, 1, 2, 5, 10, 30} (seconds).

\subsection{Approval Metrics}

\begin{table}[h]
\centering
\begin{tabular}{|p{4.5cm}|p{1.5cm}|p{3.5cm}|p{5.5cm}|}
\hline
\textbf{Metric} & \textbf{Type} & \textbf{Labels} & \textbf{Description} \\
\hline
\texttt{gert\_approval\_wait\_seconds} & Histogram & \texttt{step\_id} & Time from approval request to decision (approved or rejected) \\
\hline
\texttt{gert\_approvals\_total} & Counter & \texttt{decision}, \texttt{step\_id} & Total approval decisions (\texttt{approved}, \texttt{rejected}) \\
\hline
\end{tabular}
\caption{Approval metrics}
\end{table}

\textbf{Histogram Buckets:} \texttt{gert\_approval\_wait\_seconds} uses buckets: \texttt{10, 30, 60, 300, 600, 1800, 3600, 7200, 14400, 28800, 86400} (seconds, up to 24 hours).

\subsection{Extension Metrics}

\begin{table}[h]
\centering
\begin{tabular}{|p{4.5cm}|p{1.5cm}|p{3.5cm}|p{5.5cm}|}
\hline
\textbf{Metric} & \textbf{Type} & \textbf{Labels} & \textbf{Description} \\
\hline
\texttt{gert\_extension\_crashes\_total} & Counter & \texttt{extension\_id} & Extension process crashes \\
\hline
\texttt{gert\_extension\_loaded} & Gauge & \texttt{extension\_id} & 1 if extension is loaded, 0 otherwise \\
\hline
\end{tabular}
\caption{Extension metrics}
\end{table}

\subsection{Metrics Endpoint}

When running in server mode, metrics are exposed at:

\begin{verbatim}
GET http://localhost:8080/metrics
\end{verbatim}

Configure the server:

\begin{minted}{bash}
gert serve --http --addr=:8080
\end{minted}

The response is in Prometheus text exposition format. Example:

\begin{verbatim}
# TYPE gert_runs_total counter
gert_runs_total{outcome="success",runbook_id="deploy-app-v3"} 42
gert_runs_total{outcome="failure",runbook_id="deploy-app-v3"} 3

# TYPE gert_run_duration_seconds histogram
gert_run_duration_seconds_bucket{outcome="success",runbook_id="deploy-app-v3",le="1"} 5
gert_run_duration_seconds_bucket{outcome="success",runbook_id="deploy-app-v3",le="5"} 18
gert_run_duration_seconds_bucket{outcome="success",runbook_id="deploy-app-v3",le="+Inf"} 42
gert_run_duration_seconds_sum{outcome="success",runbook_id="deploy-app-v3"} 187.4
gert_run_duration_seconds_count{outcome="success",runbook_id="deploy-app-v3"} 42

# TYPE gert_active_runs gauge
gert_active_runs 3
\end{verbatim}

\subsection{Push-Based Metrics (OTLP)}

For environments where scraping is not feasible (e.g., ephemeral CI jobs), use OTLP metrics export:

\begin{minted}{bash}
export OTEL_EXPORTER_OTLP_ENDPOINT=http://collector:4317
export OTEL_METRICS_EXPORTER=otlp
gert run deploy.yaml
\end{minted}

The runtime pushes metrics to the OTLP endpoint at configurable intervals (default: 60 seconds).

% ================================================================
\section{Structured Logging}
% ================================================================

Gert emits structured JSON logs to \texttt{stderr}. Every log line is a complete JSON object.

\subsection{Log Line Format}

\begin{minted}{json}
{
  "timestamp": "2026-04-18T14:30:22.123456Z",
  "level": "info",
  "msg": "step execution started",
  "run_id": "20260418T143022-a3f9c1",
  "step_id": "backup-db",
  "step_type": "cli",
  "actor": "alice@example.com"
}
\end{minted}

\subsection{Required Fields}

\begin{itemize}
  \item \texttt{timestamp} — RFC 3339 format~\cite{rfc3339} with microsecond precision, UTC timezone
  \item \texttt{level} — one of: \texttt{debug}, \texttt{info}, \texttt{warn}, \texttt{error}
  \item \texttt{msg} — human-readable message describing the event
\end{itemize}

\subsection{Contextual Fields (When Applicable)}

\begin{itemize}
  \item \texttt{run\_id} — run identifier (present in all logs within run context)
  \item \texttt{step\_id} — step identifier (present in all logs within step context)
  \item \texttt{step\_type} — step type (\texttt{cli}, \texttt{manual}, \texttt{tool}, \texttt{invoke})
  \item \texttt{actor} — user or service principal executing the run
  \item \texttt{tool\_name} — tool identifier (for tool invocation logs)
  \item \texttt{error} — error message (for \texttt{error} level logs)
  \item \texttt{exit\_code} — process exit code (for CLI step completions)
\end{itemize}

\subsection{Log Levels}

\begin{itemize}
  \item \textbf{debug} — Detailed diagnostic information (e.g., input resolution, policy evaluation details)
  \item \textbf{info} — Operational milestones (e.g., step started, run completed)
  \item \textbf{warn} — Recoverable issues (e.g., retry attempt, fallback to default)
  \item \textbf{error} — Failures requiring attention (e.g., step failure, governance denial)
\end{itemize}

Configure via environment variable or CLI flag:

\begin{minted}{bash}
export GERT_LOG_LEVEL=debug
gert run deploy.yaml
\end{minted}

Or:

\begin{minted}{bash}
gert run deploy.yaml --log-level=debug
\end{minted}

Default: \texttt{info}.

\subsection{Sensitive Data Handling}

Log lines MUST NOT contain sensitive input values. The runtime applies the same redaction rules as the trace file (see \S~12.4).

\textbf{Safe:}
\begin{minted}{json}
{"msg": "resolved input 'db_password' from provider 'vault'"}
\end{minted}

\textbf{Unsafe:}
\begin{minted}{json}
{"msg": "resolved input 'db_password' = 'hunter2'"}
\end{minted}

When logging errors that may contain sensitive data, scrub or omit the error message:

\begin{minted}{json}
{"msg": "input resolution failed", "input_name": "db_password", 
 "provider": "vault", "error": "connection timeout"}
\end{minted}

\subsection{Log Output Configuration}

By default, logs are written to \texttt{stderr}. To write to a file:

\begin{minted}{bash}
gert run deploy.yaml --log-output=file:/var/log/gert/gert.log
\end{minted}

This writes logs to the file \textit{in addition to} \texttt{stderr}.

For syslog output:

\begin{minted}{bash}
gert run deploy.yaml --log-output=syslog://localhost:514
\end{minted}

\subsection{Log Aggregation and Correlation}

Every log line in run context includes \texttt{run\_id}, enabling log aggregation by run:

\textbf{Loki Query:}
\begin{verbatim}
{job="gert"} | json | run_id="20260418T143022-a3f9c1"
\end{verbatim}

\textbf{Elasticsearch Query:}
\begin{minted}{go}
GET /gert-logs-*/_search
{
  "query": {
    "term": {"run_id": "20260418T143022-a3f9c1"}
  }
}
\end{minted}

This returns all log lines for the specified run, ordered by timestamp.

% ================================================================
\section{Health and Readiness Endpoints}
% ================================================================

When running in server mode (\texttt{gert serve --http}), the runtime exposes HTTP endpoints for health probes. These are designed for Kubernetes liveness and readiness probes.

\subsection{Liveness Endpoint}

\begin{verbatim}
GET /health
\end{verbatim}

\textbf{Purpose:} Indicates whether the gert process is alive.

\textbf{Response (200 OK):}
\begin{minted}{json}
{
  "status": "ok",
  "version": ".0",
  "uptime_seconds": 3600
}
\end{minted}

\textbf{Use Case:} Kubernetes liveness probe to restart crashed pods.

\subsection{Readiness Endpoint}

\begin{verbatim}
GET /ready
\end{verbatim}

\textbf{Purpose:} Indicates whether the runtime is ready to accept new runs.

\textbf{Response (200 OK):}
\begin{minted}{json}
{
  "status": "ready",
  "extensions_loaded": 3,
  "extensions_failed": 0
}
\end{minted}

\textbf{Response (503 Service Unavailable):}
\begin{minted}{json}
{
  "status": "not_ready",
  "reason": "extension 'aws-tools' failed to load",
  "extensions_loaded": 2,
  "extensions_failed": 1
}
\end{minted}

\textbf{Use Case:} Kubernetes readiness probe to delay traffic until extensions are loaded.

\subsection{Kubernetes Probe Configuration}

\begin{minted}{yaml}
apiVersion: v1
kind: Pod
metadata:
  name: gert-server
spec:
  containers:
  - name: gert
    image: gert:.0
    command: ["gert", "serve", "--http", "--addr=:8080"]
    livenessProbe:
      httpGet:
        path: /health
        port: 8080
      initialDelaySeconds: 5
      periodSeconds: 10
    readinessProbe:
      httpGet:
        path: /ready
        port: 8080
      initialDelaySeconds: 10
      periodSeconds: 5
\end{minted}

% ================================================================
\section{Diagnostics Commands}
% ================================================================

Gert provides CLI commands for self-diagnosis and trace inspection.

\subsection{\texttt{gert diagnose}}

\textbf{Purpose:} Run pre-flight checks to verify that the runtime environment is correctly configured.

\textbf{Usage:}
\begin{minted}{bash}
gert diagnose [--runbook=<path>]
\end{minted}

\textbf{Checks Performed:}
\begin{enumerate}
  \item \textbf{Runbook Schema Validation:} Parse and validate the runbook YAML against JSON Schema
  \item \textbf{Tool Availability:} Verify that all declared tools (\texttt{.tool.yaml}) are executable and findable in \texttt{PATH} or \texttt{.runbook/tools/}
  \item \textbf{Input Provider Connectivity:} Test that all input providers (\texttt{.provider.yaml}) can be reached (e.g., Vault is reachable, API keys are valid)
  \item \textbf{Extension Loading:} Attempt to load all declared extensions and report any load failures
  \item \textbf{Governance Policy Syntax:} Validate OPA policies (if policy-as-code is configured)
\end{enumerate}

\textbf{Output (Success):}
\begin{verbatim}
✓ Runbook schema validation passed
✓ Tool 'kubectl' found at /usr/local/bin/kubectl
✓ Tool 'aws' found at /usr/local/bin/aws
✓ Input provider 'vault' reachable at https://vault:8200
✓ Extension 'aws-tools' loaded successfully
✓ Extension 'k8s-tools' loaded successfully

Diagnostics: PASS
\end{verbatim}

\textbf{Output (Failure):}
\begin{verbatim}
✓ Runbook schema validation passed
✗ Tool 'kubectl' not found in PATH or .runbook/tools/
✓ Tool 'aws' found at /usr/local/bin/aws
✗ Input provider 'vault' unreachable: connection refused
✓ Extension 'aws-tools' loaded successfully
✗ Extension 'k8s-tools' failed to load: missing dependency

Diagnostics: FAIL (3 errors)
\end{verbatim}

\textbf{Exit Code:} 0 if all checks pass, 1 if any check fails.

\subsection{\texttt{gert trace show}}

\textbf{Purpose:} Pretty-print a completed run's trace from the JSONL file.

\textbf{Usage:}
\begin{minted}{bash}
gert trace show <run_id>
\end{minted}

\textbf{Output:}
\begin{verbatim}
Run: 20260418T143022-a3f9c1
Runbook: deploy-app-v3
Actor: alice@example.com
Started: 2026-04-18T14:30:22Z
Completed: 2026-04-18T14:33:15Z
Duration: 2m53s
Outcome: success

Steps:
  [0] backup-db (cli) - 12.3s - success
    └─ tool: ssh - 12.1s - exit_code=0
  [1] deploy (tool) - 45.2s - success
    └─ tool: kubectl - 44.8s - exit_code=0
  [2] verify (manual) - 2m10s - success
    └─ approval: alice@example.com - granted
\end{verbatim}

\subsection{\texttt{gert trace export}}

\textbf{Purpose:} Convert a JSONL trace to OpenTelemetry OTLP format for import into Jaeger or Zipkin.

\textbf{Usage:}
\begin{minted}{bash}
gert trace export <run_id> --format=otlp --output=trace.json
\end{minted}

\textbf{Supported Formats:}
\begin{itemize}
  \item \texttt{otlp} — OTLP JSON format
  \item \texttt{jaeger} — Jaeger JSON format
  \item \texttt{zipkin} — Zipkin JSON format
\end{itemize}

\textbf{Use Case:} Retrospectively analyze a run in a trace visualization tool when the runtime was not configured to export traces at execution time.

\subsection{\texttt{gert run list}}

\textbf{Purpose:} List recent runs with status.

\textbf{Usage:}
\begin{minted}{bash}
gert run list [--limit=N] [--json]
\end{minted}

\textbf{Output (Human-Readable):}
\begin{verbatim}
RUN_ID                       RUNBOOK          ACTOR                OUTCOME    DURATION
20260418T143022-a3f9c1       deploy-app-v3    alice@example.com    success    2m53s
20260418T140015-f3a7b2       rollback-prod      bob@example.com      failure    45s
20260418T135512-c9e8d1       backup-db        cron@system          success    3m12s
\end{verbatim}

\textbf{Output (JSON, for CI integration):}
\begin{minted}{json}
[
  {
    "run_id": "20260418T143022-a3f9c1",
    "runbook_id": "deploy-app-v3",
    "actor": "alice@example.com",
    "outcome": "success",
    "duration_ms": 173000,
    "started_at": "2026-04-18T14:30:22Z",
    "completed_at": "2026-04-18T14:33:15Z"
  },
  ...
]
\end{minted}

\textbf{Use Case:} CI pipelines can parse JSON output to determine run status.

% ================================================================
\section{Integration Recipes}
% ================================================================

This section provides practical integration patterns for common observability platforms.

\subsection{Prometheus + Grafana}

\textbf{Step 1: Start gert server}
\begin{minted}{bash}
gert serve --http --addr=:8080
\end{minted}

\textbf{Step 2: Configure Prometheus scrape target}

Add to \texttt{prometheus.yml}:
\begin{minted}{yaml}
scrape_configs:
  - job_name: 'gert'
    static_configs:
      - targets: ['localhost:8080']
    metrics_path: '/metrics'
    scrape_interval: 15s
\end{minted}

\textbf{Step 3: Import Grafana dashboard}

Gert provides a reference Grafana dashboard JSON at \texttt{.runbook/observability/grafana-dashboard.json}. Import it into Grafana to visualize:
\begin{itemize}
  \item Run throughput (runs/minute)
  \item Success vs. failure rate
  \item p50, p95, p99 run duration
  \item Active runs (gauge)
  \item Step execution time by type
  \item Tool invocation latency
\end{itemize}

\subsection{Jaeger (Distributed Tracing)}

\textbf{Step 1: Start Jaeger (all-in-one)}
\begin{minted}{bash}
docker run -d --name jaeger \
  -p 16686:16686 \
  -p 4317:4317 \
  jaegertracing/all-in-one:latest
\end{minted}

\textbf{Step 2: Configure gert to export traces}
\begin{minted}{bash}
export OTEL_EXPORTER_OTLP_ENDPOINT=http://localhost:4317
gert run deploy.yaml
\end{minted}

\textbf{Step 3: View traces in Jaeger UI}

Open \url{http://localhost:16686}, select service \texttt{gert}, and search for traces.

\subsection{Loki + Grafana (Log Aggregation)}

\textbf{Step 1: Configure Promtail to scrape gert logs}

Add to \texttt{promtail-config.yaml}:
\begin{minted}{yaml}
scrape_configs:
  - job_name: gert
    static_configs:
      - targets:
          - localhost
        labels:
          job: gert
          __path__: /var/log/gert/*.log
    pipeline_stages:
      - json:
          expressions:
            timestamp: timestamp
            level: level
            msg: msg
            run_id: run_id
      - labels:
          level:
          run_id:
\end{minted}

\textbf{Step 2: Start gert with file logging}
\begin{minted}{bash}
gert run deploy.yaml --log-output=file:/var/log/gert/gert.log
\end{minted}

\textbf{Step 3: Query logs in Grafana Explore}
\begin{verbatim}
{job="gert"} | json | run_id="20260418T143022-a3f9c1"
\end{verbatim}

\subsection{Datadog / New Relic (OTLP-Compatible SaaS)}

Both platforms support OpenTelemetry OTLP ingestion.

\textbf{Datadog:}
\begin{minted}{bash}
export OTEL_EXPORTER_OTLP_ENDPOINT=https://api.datadoghq.com
export DD_API_KEY=<your-api-key>
gert run deploy.yaml
\end{minted}

\textbf{New Relic:}
\begin{minted}{bash}
export OTEL_EXPORTER_OTLP_ENDPOINT=https://otlp.nr-data.net:4317
export NEW_RELIC_API_KEY=<your-api-key>
gert run deploy.yaml
\end{minted}

Consult platform documentation for exact endpoint URLs and authentication headers.

% ================================================================
\section{Observability for CI/CD Pipelines}
% ================================================================

Gert provides machine-readable output for CI/CD integration.

\subsection{JSON Output Mode}

\begin{minted}{bash}
gert run deploy.yaml --output=json
\end{minted}

On completion, emits a JSON summary to \texttt{stdout}:

\begin{minted}{json}
{
  "run_id": "20260418T143022-a3f9c1",
  "outcome": "success",
  "step_count": 3,
  "failed_step": null,
  "duration_ms": 173000,
  "started_at": "2026-04-18T14:30:22Z",
  "completed_at": "2026-04-18T14:33:15Z"
}
\end{minted}

\textbf{Exit Code:}
\begin{itemize}
  \item 0 if \texttt{outcome == "success"}
  \item 1 if \texttt{outcome == "failure"}
  \item 2 if \texttt{outcome == "governance-blocked"}
  \item 3 if \texttt{outcome == "cancelled"}
\end{itemize}

The run summary object above is \textbf{unchanged} by the enum constraint
MVP (design/gert/sections/03-schema-vnext.tex \S\ref{para:enum-constraint}):
it carries no enum metadata. An enum violation surfaces as an
\texttt{ENUM-007}/\texttt{ENUM-008}/\texttt{ENUM-009} error object using the
existing structured error envelope (design/gert/sections/03d-parse-time-enforcement.tex
\S Error Message Format), subject to the redaction rules of
design/gert/sections/08-security-and-trust.tex \S Secret Resolution and
Redaction for any declaration carrying a sensitivity marking.

\subsection{GitHub Actions Example}

\begin{minted}{yaml}
- name: Run gert runbook
  id: gert_run
  run: |
    OUTPUT=$(gert run deploy.yaml --output=json)
    echo "$OUTPUT" | jq .
    RUN_ID=$(echo "$OUTPUT" | jq -r .run_id)
    OUTCOME=$(echo "$OUTPUT" | jq -r .outcome)
    echo "run_id=$RUN_ID" >> $GITHUB_OUTPUT
    echo "outcome=$OUTCOME" >> $GITHUB_OUTPUT

- name: Upload trace artifact
  if: always()
  uses: actions/upload-artifact@v3
  with:
    name: gert-trace
    path: .runbook/runs/${{ steps.gert_run.outputs.run_id }}/trace.jsonl
\end{minted}

\subsection{GitLab CI Example}

\begin{minted}{yaml}
deploy:
  script:
    - gert run deploy.yaml --output=json | tee run-output.json
    - RUN_ID=$(jq -r .run_id run-output.json)
    - echo "Run ID: $RUN_ID"
  artifacts:
    when: always
    paths:
      - .runbook/runs/
    reports:
      junit: .runbook/runs/*/trace.jsonl
\end{minted}

% ================================================================
\section{Summary}
% ================================================================

This chapter specifies a production-grade observability model for gert grounded in industry standards:

\begin{itemize}
  \item \textbf{Three-pillar model:} Metrics (Prometheus), distributed tracing (OpenTelemetry), structured logging (JSON)
  \item \textbf{Span hierarchy:} \texttt{gert.run} → \texttt{gert.step} → \texttt{gert.tool.invoke / gert.governance.check / gert.input.resolve}
  \item \textbf{Context propagation:} W3C Trace Context headers and OpenTelemetry baggage for cross-system correlation
  \item \textbf{Metrics catalog:} 9 metrics covering runs, steps, tools, approvals, and extensions
  \item \textbf{Health endpoints:} \texttt{/health} (liveness), \texttt{/ready} (readiness), \texttt{/metrics} (Prometheus)
  \item \textbf{Diagnostics commands:} \texttt{gert diagnose}, \texttt{gert trace show}, \texttt{gert trace export}, \texttt{gert run list}
  \item \textbf{Integration recipes:} Prometheus, Grafana, Jaeger, Loki, Datadog, New Relic
  \item \textbf{CI/CD support:} JSON output mode with exit codes for pipeline integration
\end{itemize}

Operators gain visibility into run health, performance bottlenecks, and failure root causes. The design ensures portability across observability platforms by adhering to open standards~\cite{otel-spec, w3c-trace-context, google-sre-book}.
